\section{Results (paper scaffolding)}

This section defines stable figure slots referenced by the paper.
Figures are generated by \texttt{python scripts/make\_paper.py}.

\subsection{Descriptive bubble statistics}
\begin{figure}[ht]
  \centering
  \includegraphics[width=0.9\textwidth]{fig\_bubble\_descriptives.png}
  \caption{Descriptive bubble diagnostics (placeholder caption).}
\end{figure}

\subsection{Bubble overlap chart (Gantt-style)}
\begin{figure}[ht]
  \centering
  \includegraphics[width=0.95\textwidth]{fig\_overlap\_gantt.png}
  \caption{Bubble overlap chart across firms (placeholder caption).}
\end{figure}

\subsection{Aggregate overlap network}
\begin{figure}[ht]
  \centering
  \includegraphics[width=0.85\textwidth]{fig\_overlap\_network.png}
  \caption{Aggregate bubble-overlap network visualization (directed, overlap-weighted).}
\end{figure}

\begin{table}[ht]
\centering
\small
\begin{tabular}{l r r r r}
\hline
Network & Nodes & Edges & Density & Weak CCs \\
\hline
Aggregate & 19 & 246 & 0.7193 & 1 \\
Temporal (avg) & 19.0 & 12.9 & 0.0377 & 15.8 \\
\hline
\end{tabular}
\caption{Baseline network diagnostics (aggregate and temporal average).}
\label{tab:network_summary}
\end{table}

\noindent\textbf{Top nodes (aggregate)}\\
\begin{verbatim}
in-degree: LPO:18; SNP:18; TEL:18; BRD:17; SNG:17; TGN:17; TLV:17; FNC:14; FP:14; PUR:14
out-degree: FNC:18; ATB:16; TTS:16; FP:15; PUR:15; SNN:15; TLV:15; AP1:14; MED:14; TRP:14
in-weighted-degree: FP:3.21e+03; TLV:3.14e+03; AP1:2.94e+03; TEL:2.67e+03; SNP:2.65e+03; LPO:2.53e+03; DIG:2.18e+03; PUR:1.99e+03; BRD:1.93e+03; MED:1.47e+03
out-weighted-degree: FNC:5.34e+03; ATB:4.15e+03; TTS:3.65e+03; PUR:2.8e+03; DIG:2.3e+03; FP:2.13e+03; TLV:1.98e+03; TRP:1.77e+03; MED:1.76e+03; SDP:1.47e+03
\end{verbatim}

\noindent\textbf{Top nodes (temporal average)}\\
\begin{verbatim}
in-degree(avg): FP:1.23; TLV:1.21; AP1:1.13; TEL:1.07; SNP:1.02; LPO:0.985; DIG:0.838; PUR:0.78; BRD:0.758; TGN:0.564
out-degree(avg): FNC:2.04; ATB:1.6; TTS:1.41; PUR:1.08; DIG:0.888; FP:0.828; TLV:0.768; MED:0.689; TRP:0.688; SDP:0.573
weighted-degree(avg): FNC:2.18; FP:2.06; ATB:2.02; TLV:1.98; PUR:1.86; TTS:1.8; DIG:1.73; AP1:1.64; LPO:1.36; SNP:1.31
\end{verbatim}


\begin{figure}[ht]
  \centering
  \includegraphics[width=0.95\textwidth]{fig\_degree\_distributions.png}
  \caption{Degree distributions for the aggregate overlap network (simple histogram diagnostics).}
\end{figure}

\subsection{Temporal centrality dynamics}
\begin{figure}[ht]
  \centering
  \includegraphics[width=0.95\textwidth]{fig\_centrality\_dynamics.png}
  \caption{Centrality dynamics over time for the top nodes (by average eigenvector centrality).}
\end{figure}

\begin{figure}[ht]
  \centering
  \includegraphics[width=0.95\textwidth]{fig\_centrality\_heatmap.png}
  \caption{Eigenvector centrality heatmap across firms (temporal snapshots).}
\end{figure}

\IfFileExists{tables/table_centrality_summary.tex}{
  \begin{table}[ht]
\centering
\small
\begin{tabular}{l r r}
\hline
Company & Mean eigenvector & Std eigenvector \\
\hline
SNN & 0.2272 & 0.2170 \\
MED & 0.1985 & 0.1943 \\
TRP & 0.1949 & 0.2091 \\
FP & 0.1888 & 0.1640 \\
FNC & 0.1882 & 0.1678 \\
PUR & 0.1722 & 0.1639 \\
AP1 & 0.1721 & 0.1789 \\
ATB & 0.1608 & 0.1483 \\
TLV & 0.1589 & 0.1457 \\
TTS & 0.1558 & 0.1547 \\
\hline
\end{tabular}
\caption{Top-10 nodes by average eigenvector centrality (temporal).}
\label{tab:centrality_top10}
\end{table}

}{
  \noindent\textit{Centrality table not available (run make\_paper without \texttt{--skip-centrality}).}
}

\begin{figure}[ht]
  \centering
  \includegraphics[width=0.95\textwidth]{fig\_centrality\_top\_nodes.png}
  \caption{Eigenvector centrality time series for the top-5 nodes (average over time).}
\end{figure}

\subsection{TGNN performance (optional)}
\begin{figure}[ht]
  \centering
  \includegraphics[width=0.75\textwidth]{fig\_tgnn\_performance.png}
  \caption{TGNN forecast performance (optional). If TGNN is not executed, this panel is an auto-generated placeholder.}
\end{figure}

\subsection{FRM network properties}

\IfFileExists{tables/table_frm_network_summary.tex}{
  \begin{table}[ht]
\centering
\small
\begin{tabular}{l r r r r}
\hline
Network & Nodes & Edges & Density & Weak CCs \\
\hline
FRM (agg) & 20 & 116 & 0.3053 & 9 \\
FRM (avg) & 20.0 & 33.0 & 0.0868 & 14.2 \\
\hline
\end{tabular}
\caption{FRM network diagnostics (aggregate and temporal average).}
\label{tab:frm_network_summary}
\end{table}

\noindent\textbf{Top nodes (aggregate)}\\
\begin{verbatim}
in-degree: BRD:11; SNG:11; SNN:11; TLV:11; DIG:10; FP:10; LPO:10; SNP:10; TGN:10; TRP:9
out-degree: BRD:11; SNG:11; SNN:11; TLV:11; DIG:10; FP:10; LPO:10; SNP:10; TGN:10; TRP:9
in-weighted-degree: SNN:86.6; BRD:85.5; TLV:75.5; SNP:44.5; FP:38.4; SNG:35.3; TGN:33.2; LPO:31.6; DIG:24.3; TEL:19.6
out-weighted-degree: TGN:82.8; SNG:67.6; TLV:64.7; BRD:55.3; SNP:52.9; FP:36.4; SNN:35.5; TEL:34.7; LPO:24; DIG:21.1
\end{verbatim}

\noindent\textbf{Top nodes (temporal average)}\\
\begin{verbatim}
in-degree(avg): BRD:4.72; TLV:4.72; SNN:3.9; SNG:3.62; SNP:3.52; TGN:3.17; FP:2.62; DIG:2; LPO:1.88; TEL:1.57
out-degree(avg): TLV:4.72; BRD:4.28; SNG:4.22; TGN:3.73; SNP:3.58; SNN:3.42; FP:2.65; LPO:2.02; TEL:1.65; DIG:1.62
weighted-degree(avg): BRD:3.52; TLV:3.5; SNN:3.05; TGN:2.9; SNG:2.57; SNP:2.44; FP:1.87; LPO:1.39; TEL:1.36; DIG:1.14
\end{verbatim}

}{
  \noindent\textit{FRM table not available (run make\_paper with \texttt{--run-frm}).}
}

\IfFileExists{figures/fig_frm_degree_distributions.png}{
  \begin{figure}[ht]
    \centering
    \includegraphics[width=0.95\textwidth]{fig\_frm\_degree\_distributions.png}
    \caption{Degree distributions for the aggregate FRM network (simple histogram diagnostics).}
  \end{figure}
}{}

\subsection{Bubble vs FRM comparison}
\IfFileExists{figures/fig_frm_overlap_vs_bubble.png}{
  \begin{figure}[ht]
    \centering
    \includegraphics[width=0.75\textwidth]{fig\_frm\_overlap\_vs\_bubble.png}
    \caption{Similarity between the aggregate bubble-overlap and FRM networks using simple, interpretable metrics.}
  \end{figure}
}{}

\subsection{Bubble vs FRM similarity over time}

We also quantify similarity between the two systemic-risk lenses at the snapshot level, not only in aggregate.
High similarity indicates that the two constructions highlight the same set of central firms and concentrated edge structure, while low similarity suggests that they capture different channels (bubble-timing overlap versus tail-dependence regression links).

\IfFileExists{tables/table_bubble_vs_frm_similarity.tex}{
  \begin{table}[ht]
\centering
\small
\begin{tabular}{l r r r r}
\hline
Metric & Mean & Std & Min & Max \\
\hline
spearman\_eigenvector\_rank & 0.042 & 0.210 & -0.609 & 0.599 \\
kendall\_eigenvector\_rank & 0.031 & 0.151 & -0.481 & 0.450 \\
top5\_jaccard\_nodes & 0.191 & 0.099 & 0.000 & 0.667 \\
top10\_jaccard\_nodes & 0.389 & 0.104 & 0.111 & 0.667 \\
top\_edges\_jaccard & 0.023 & 0.036 & 0.000 & 0.179 \\
top\_edges\_weight\_corr & NA & NA & NA & NA \\
\hline
\end{tabular}
\caption{Bubble vs FRM similarity over time (aligned by nearest prior date).}
\label{tab:bubble_vs_frm_similarity}
\end{table}

}{
  \noindent\textit{Bubble-vs-FRM similarity table not available (run make\_paper with \texttt{--run-compare} and \texttt{--run-frm}).}
}

\IfFileExists{figures/fig_bubble_vs_frm_similarity_timeseries.png}{
  \begin{figure}[ht]
    \centering
    \includegraphics[width=0.95\textwidth]{fig\_bubble\_vs\_frm\_similarity\_timeseries.png}
    \caption{Time-aware Bubble vs FRM similarity (aligned by nearest prior FRM date). \texttt{minimal} builds are for smoke-tests/CI and may restrict FRM (e.g., $N=20$ firms, shorter window); \texttt{full} builds should be used for final reported results.}
  \end{figure}
}{}
